\appendix
\clearpage
\markboth{}{}
\pagestyle{fancy}
\addcontentsline{toc}{part}{ХАВСРАЛТ}



% Хавсралтын нэр. Хавсралт гэдэг үг агуулахгүй
\chapter{Кодын хэрэгжүүлэлт}


\section{\textit{\textbf{ApplicationContext.java  }}}
\begin{lstlisting}[language=Java]
package mn.edu.num.container;

import mn.edu.num.exception.BeanCreationException;
import mn.edu.num.exception.CircularDependencyException;
import mn.edu.num.exception.NoSuchBeanException;
import mn.edu.num.scanner.ClassPathScanner;

import java.util.*;

public class ApplicationContext implements DependencyInjector.ApplicationContextRef {

    private final BeanRegistry registry = new BeanRegistry();
    private final DependencyInjector injector = new DependencyInjector(registry);
    // Circular dependency илрүүлэх
    private final Set<String> inCreation = new HashSet<>();

    public ApplicationContext(String basePackage) {
        System.out.println("[Context] Эхэлж байна package:" + basePackage);
        ClassPathScanner scanner = new ClassPathScanner(basePackage);
        List<BeanDefinition> definitions = scanner.scan();

        // Бүртгэх
        for (BeanDefinition def : definitions) {
            registry.register(def);
        }

        // Singleton bean-уудыг урьдчилан үүсгэх
        for (BeanDefinition def : definitions) {
            if (def.getScope() == ScopeType.SINGLETON) {
                getBean(def.getBeanName());
            }
        }
        System.out.println("[Context] Container бэлэн боллоо.");
    }

    @Override
    public Object getBean(String name) {
        BeanDefinition def = registry.getDefinition(name);

        // Кэш шалгахаас ӨМНӨ тойрог хамаарал шалгана
        if (inCreation.contains(name)) {
            throw new CircularDependencyException(
                    "Тойрог хамаарал илэрлээ:" + name
                            + " | Үүсгэгдэж байгаа: " + inCreation);
        }
        if (def.getScope() == ScopeType.SINGLETON) {
            if (registry.hasSingleton(name)) {
                return registry.getSingleton(name);
            }
            return createBean(def);
        } else {
            // Prototype — шинэ instance үүсгэнэ
            return createBean(def);
        }
    }

    public <T> T getBean(Class<T> type) {
        List<BeanDefinition> candidates = registry.findByType(type);
        if (candidates.isEmpty())
            throw new NoSuchBeanException("Bean олдсонгүй:" + type.getName());
        if (candidates.size() > 1)
            throw new NoSuchBeanException("Олон bean тохирч байна:" + type.getName());
        return type.cast(getBean(candidates.get(0).getBeanName()));
    }

    private Object createBean(BeanDefinition def) {
        String name = def.getBeanName();

        // Circular dependency шалгах
        if (inCreation.contains(name)) {
            throw new CircularDependencyException(
                    "Тойрог хамаарал илэрлээ: " + name);
        }
        inCreation.add(name);

        try {
            Object instance = def.getBeanClass()
                    .getDeclaredConstructor()
                    .newInstance();
            // Singleton бол cache-д хадгална
            if (def.getScope() == ScopeType.SINGLETON) {
                registry.cacheSingleton(name, instance);
            }
            // Dependency inject хийнэ
            injector.inject(instance, def, this);
            return instance;

        } catch (CircularDependencyException e) {
            throw e;
        } catch (Exception e) {
            throw new BeanCreationException(
                    "Bean үүсгэхэд алдаа:" + name, e);
        } finally {
            inCreation.remove(name);
        }
    }
}

\end{lstlisting}

\clearpage
\section{\textit{\textbf{ClassPathScanner.java}}}
\begin{lstlisting}[language=Java]
package mn.edu.num.scanner;

import mn.edu.num.annotation.Component;
import mn.edu.num.container.BeanDefinition;
import mn.edu.num.container.ScopeType;
import mn.edu.num.annotation.Scope;
import java.io.File;
import java.net.URL;
import java.util.ArrayList;
import java.util.List;

public class ClassPathScanner {

    private final String basePackage;

    public ClassPathScanner(String basePackage) {
        this.basePackage = basePackage;
    }

    public List<BeanDefinition> scan() {
        List<BeanDefinition> definitions = new ArrayList<>();
        String path = basePackage.replace('.', '/');
        URL resource = Thread.currentThread()
                .getContextClassLoader()
                .getResource(path);

        if (resource == null) return definitions;

        File dir = new File(resource.getFile());
        scanDirectory(dir, basePackage, definitions);
        return definitions;
    }

    private void scanDirectory(File dir, String packageName,
                               List<BeanDefinition> definitions) {
        File[] files = dir.listFiles();
        if (files == null) return;

        for (File file : files) {
            if (file.isDirectory()) {
                scanDirectory(file, packageName + "." + file.getName(), definitions);
            } else if (file.getName().endsWith(".class")) {
                String className = packageName + "."
                        + file.getName().replace(".class", "");
                processClass(className, definitions);
            }
        }
    }

    private void processClass(String className, List<BeanDefinition> definitions) {
        try {
            Class<?> clazz = Class.forName(className);
            if (!clazz.isAnnotationPresent(Component.class)) return;

            Component component = clazz.getAnnotation(Component.class);
            String beanName = component.value().isEmpty()
                    ? decapitalize(clazz.getSimpleName())
                    : component.value();

            ScopeType scope = ScopeType.SINGLETON;
            if (clazz.isAnnotationPresent(Scope.class)) {
                String scopeVal = clazz.getAnnotation(Scope.class).value();
                scope = "prototype".equalsIgnoreCase(scopeVal)
                        ? ScopeType.PROTOTYPE : ScopeType.SINGLETON;
            }

            definitions.add(new BeanDefinition(beanName, clazz, scope));
            System.out.println("[Scanner] Илэрлээ:" + beanName + ":" + clazz.getName());

        } catch (ClassNotFoundException e) {
            System.err.println("[Scanner] Класс ачаалахад алдаа:" + className);
        }
    }

    private String decapitalize(String name) {
        return Character.toLowerCase(name.charAt(0)) + name.substring(1);
    }
}
\end{lstlisting}


\section{\textit{\textbf{DependencyInjector.java}}}
\begin{lstlisting}[language=Java]
package mn.edu.num.container;

import mn.edu.num.annotation.Autowired;
import mn.edu.num.annotation.Qualifier;
import mn.edu.num.exception.BeanCreationException;
import mn.edu.num.exception.NoSuchBeanException;

import java.lang.reflect.Field;
import java.util.List;

public class DependencyInjector {

    private final BeanRegistry registry;

    public DependencyInjector(BeanRegistry registry) {
        this.registry = registry;
    }

    public void inject(Object instance, BeanDefinition definition,
                       ApplicationContextRef contextRef) {
        for (Field field : definition.getBeanClass().getDeclaredFields()) {
            if (!field.isAnnotationPresent(Autowired.class)) continue;

            Class<?> fieldType = field.getType();
            String beanName = null;

            // @Qualifier байвал нэрээр хайна
            if (field.isAnnotationPresent(Qualifier.class)) {
                beanName = field.getAnnotation(Qualifier.class).value();
            }

            Object dependency;
            if (beanName != null) {
                dependency = contextRef.getBean(beanName);
            } else {
                // Төрлөөр хайна
                List<BeanDefinition> candidates = registry.findByType(fieldType);
                if (candidates.isEmpty()) {
                    throw new NoSuchBeanException(
                            "Dependency олдсонгүй:" + fieldType.getName()
                                    + " [" + definition.getBeanName() + "]");
                }
                if (candidates.size() > 1) {
                    throw new NoSuchBeanException(
                            "Олон bean тохирч байна:" + fieldType.getName()
                                    + ".@Qualifier ашиглана уу.");
                }
                dependency = contextRef.getBean(candidates.get(0).getBeanName());
            }

            try {
                field.setAccessible(true);
                field.set(instance, dependency);
                System.out.println("[Injector] Inject:"
                        + definition.getBeanName() + "." + field.getName());
            } catch (IllegalAccessException e) {
                throw new BeanCreationException(
                        "Field inject хийхэд алдаа:" + field.getName(), e);
            }
        }
    }

    // ApplicationContext-тэй circular reference үүсгэхгүйн тулд interface ашиглана
    public interface ApplicationContextRef {
        Object getBean(String name);
    }
}
\end{lstlisting}

\section{\textit{\textbf{BeanRegistry.java}}}
\begin{lstlisting}[language=Java]
package mn.edu.num.container;

import mn.edu.num.exception.NoSuchBeanException;

import java.util.*;

public class BeanRegistry {
    // bean нэр : BeanDefinition
    private final Map<String, BeanDefinition> definitions = new HashMap<>();
    // singleton cache
    private final Map<String, Object> singletonCache = new HashMap<>();

    public void register(BeanDefinition definition) {
        definitions.put(definition.getBeanName(), definition);
    }

    public BeanDefinition getDefinition(String name) {
        BeanDefinition def = definitions.get(name);
        if (def == null) throw new NoSuchBeanException("Bean олдсонгүй:" + name);
        return def;
    }

    public Collection<BeanDefinition> getAllDefinitions() {
        return definitions.values();
    }

    public void cacheSingleton(String name, Object instance) {
        singletonCache.put(name, instance);
    }

    public Object getSingleton(String name) {
        return singletonCache.get(name);
    }

    public boolean hasSingleton(String name) {
        return singletonCache.containsKey(name);
    }

    // Төрлөөр хайх
    public List<BeanDefinition> findByType(Class<?> type) {
        List<BeanDefinition> result = new ArrayList<>();
        for (BeanDefinition def : definitions.values()) {
            if (type.isAssignableFrom(def.getBeanClass())) {
                result.add(def);
            }
        }
        return result;
    }
}

\end{lstlisting}


\section{\textit{\textbf{ContextIntegrationTest.java}}}
\begin{lstlisting}[language=Java]
import mn.edu.num.app.OrderService;
import mn.edu.num.app.UserService;
import mn.edu.num.app.qualifer.presistence.MongoUserRepository;
import mn.edu.num.app.qualifer.service.ServiceFactory;
import mn.edu.num.container.ApplicationContext;
import mn.edu.num.exception.BeanCreationException;
import mn.edu.num.exception.CircularDependencyException;
import mn.edu.num.exception.NoSuchBeanException;
import org.junit.jupiter.api.Test;

import static org.junit.jupiter.api.Assertions.*;

class ContextIntegrationTest {
    @Test
    void testDependencyInjectionEmailService() {
        ApplicationContext ctx = new ApplicationContext("mn.edu.num.app");
        UserService us = ctx.getBean(UserService.class);
        assertNotNull(us);
        assertEquals("Имэйл илгээлээ", us.register());
    }

    @Test
    void testDependencyInjectionOrderService() {
        ApplicationContext ctx = new ApplicationContext("mn.edu.num.app");
        UserService us = ctx.getBean(UserService.class);
        assertNotNull(us);
        assertEquals("Захиалга хийгдлээ!", us.orderRegister());
    }

    @Test
    void testCircularDependency() {
        CircularDependencyException ex = assertThrows(
                CircularDependencyException.class,
                () -> new ApplicationContext("mn.edu.num.circular")  // ← засав
        );
        System.out.println("Exception:" + ex.getMessage());
        assertTrue(ex.getMessage().contains("Тойрог хамаарал илэрлээ: "));
    }

    @Test
    void testNoSuchBean_ByName() {
        ApplicationContext ctx = new ApplicationContext("mn.edu.num.app");

        NoSuchBeanException ex = assertThrows(
                NoSuchBeanException.class,
                () -> ctx.getBean("байхгүйСервис")
        );

        System.out.println("Exception:" + ex.getMessage());
        assertTrue(ex.getMessage().contains("байхгүйСервис")
                || ex.getMessage().contains("Bean олдсонгүй"));
    }

    @Test
    void testNoSuchBean_ByType() {
        ApplicationContext ctx = new ApplicationContext("mn.edu.num.app");

        // String класс bean болгон бүртгэгдээгүй
        NoSuchBeanException ex = assertThrows(
                NoSuchBeanException.class,
                () -> ctx.getBean(String.class)
        );

        System.out.println("Exception:" + ex.getMessage());
        assertNotNull(ex.getMessage());
    }

    @Test
    void testBeanCreationException_NoDefaultConstructor() {

        BeanCreationException ex = assertThrows(
                BeanCreationException.class,
                () -> new ApplicationContext("mn.edu.num.broken")
                // mn.edu.num.broken дотор NoDefaultConstructor-тэй класс байна
        );

        System.out.println("Exception:" + ex.getMessage());
        assertNotNull(ex.getCause()); // Reflection-ийн үндсэн алдаа байх ёстой
    }

    @Test
    void testQualiferCheck() {
        ApplicationContext app = new ApplicationContext("mn.edu.num.app.qualifer");
        ServiceFactory sf = app.getBean(ServiceFactory.class);

        assertNotNull(sf);
        assertEquals("MongoDB Repository", sf.getRepoType());
    }

    @Test
    void testQualifer_getMongoBean_ByCustomName(){
        ApplicationContext ctx = new ApplicationContext("mn.edu.num.app.qualifer");
        Object bean = ctx.getBean("MongoRepository");

        assertNotNull(bean);
        assertInstanceOf(MongoUserRepository.class, bean);

        MongoUserRepository repository = (MongoUserRepository) bean;
        assertEquals("MongoDB Repository", repository.getType());
    }
}


\end{lstlisting}